\documentclass[a4paper,10pt]{report}\input{../0.Préambule/Préambule_cours_prof}\begin{document}

\definecolor{titlecolor}{rgb}{0.98, 0.4, 0.0}
\pagestyle{empty}
\newpage

\vspace{-2cm}
\setcounter{cmeth}{0}
\begin{meth}[troncature ou arrondi ?]

La fraction \href{https://www.numworks.com/fr/simulateur/}{$\dfrac{37}{23}$} ayant un quotient décimal infini, on souhaite ne donner une valeur approchée au millième. 

\noindent \begin{minipage}{7cm}
\begin{center}
\textbf{Troncature}
\end{center}
Pour donner la troncature au millième, on coupe le nombre décimal juste après sa troisième décimale :

\vspace{0.5cm}
\begin{center}
\huge{$1,608$} \bcstop~ \textcolor{lightgray}{$695652...$}
\end{center}

\vspace{0.9cm}
La troncature au millième de $\dfrac{37}{13}$ est $1,608$
\end{minipage}
\hspace{3mm}
\vline
\hspace{3mm}
\begin{minipage}{9cm}
\begin{center}
\textbf{Arrondi}
\end{center}
Pour donner l'arrondi au millième, il faut regarder la quatrième décimale :

\begin{enumerate}
\item[•]Si il s'agit du $0$, $1$, $2$, $3$, $4$ alors l'arrondi est équivalent à la troncature.
\item[•]Si il s'agit du $5$, $6$, $7$, $8$, $9$ alors on augment de $1$ la troisième décimale.
\end{enumerate}

\vspace{-0.5cm}
\begin{center}
\huge{$1,608$}\textcolor{red}{$6$}\textcolor{lightgray}{$695652...$}
\end{center}
La quatrième décimale étant un $6$, l'arrondi au millième de $\dfrac{37}{23}$ est donc $1,609$.
\end{minipage}
\end{meth}

\setcounter{cmeth}{0}
\begin{meth}[troncature ou arrondi ?]

La fraction \href{https://www.numworks.com/fr/simulateur/}{$\dfrac{37}{23}$} ayant un quotient décimal infini, on souhaite ne donner une valeur approchée au millième. 

\noindent \begin{minipage}{7cm}
\begin{center}
\textbf{Troncature}
\end{center}
Pour donner la troncature au millième, on coupe le nombre décimal juste après sa troisième décimale :

\vspace{0.5cm}
\begin{center}
\huge{$1,608$} \bcstop~ \textcolor{lightgray}{$695652...$}
\end{center}

\vspace{0.9cm}
La troncature au millième de $\dfrac{37}{13}$ est $1,608$
\end{minipage}
\hspace{3mm}
\vline
\hspace{3mm}
\begin{minipage}{9cm}
\begin{center}
\textbf{Arrondi}
\end{center}
Pour donner l'arrondi au millième, il faut regarder la quatrième décimale :

\begin{enumerate}
\item[•]Si il s'agit du $0$, $1$, $2$, $3$, $4$ alors l'arrondi est équivalent à la troncature.
\item[•]Si il s'agit du $5$, $6$, $7$, $8$, $9$ alors on augment de $1$ la troisième décimale.
\end{enumerate}

\vspace{-0.5cm}
\begin{center}
\huge{$1,608$}\textcolor{red}{$6$}\textcolor{lightgray}{$695652...$}
\end{center}
La quatrième décimale étant un $6$, l'arrondi au millième de $\dfrac{37}{23}$ est donc $1,609$.
\end{minipage}
\end{meth}

\setcounter{cmeth}{0}
\begin{meth}[troncature ou arrondi ?]

La fraction \href{https://www.numworks.com/fr/simulateur/}{$\dfrac{37}{23}$} ayant un quotient décimal infini, on souhaite ne donner une valeur approchée au millième. 

\noindent \begin{minipage}{7cm}
\begin{center}
\textbf{Troncature}
\end{center}
Pour donner la troncature au millième, on coupe le nombre décimal juste après sa troisième décimale :

\vspace{0.5cm}
\begin{center}
\huge{$1,608$} \bcstop~ \textcolor{lightgray}{$695652...$}
\end{center}

\vspace{0.9cm}
La troncature au millième de $\dfrac{37}{13}$ est $1,608$
\end{minipage}
\hspace{3mm}
\vline
\hspace{3mm}
\begin{minipage}{9cm}
\begin{center}
\textbf{Arrondi}
\end{center}
Pour donner l'arrondi au millième, il faut regarder la quatrième décimale :

\begin{enumerate}
\item[•]Si il s'agit du $0$, $1$, $2$, $3$, $4$ alors l'arrondi est équivalent à la troncature.
\item[•]Si il s'agit du $5$, $6$, $7$, $8$, $9$ alors on augment de $1$ la troisième décimale.
\end{enumerate}

\vspace{-0.5cm}
\begin{center}
\huge{$1,608$}\textcolor{red}{$6$}\textcolor{lightgray}{$695652...$}
\end{center}
La quatrième décimale étant un $6$, l'arrondi au millième de $\dfrac{37}{23}$ est donc $1,609$.
\end{minipage}
\end{meth}
\end{document}